\section{Research question and precise status}
\label{sec:context}

\subsection{What is being accelerated}

The input to the local problem is a finite equivalence relation encoded by
isometries between integer intervals. The integer universe may be enormous,
while the number of interval pairings and their endpoint bit lengths are
small. A connected normal surface component can be represented by an orbit
of such a relation. The immediate computational objective is to avoid
repeatedly rediscovering or transporting the same component information when
many observations share one fixed source.

An important distinction is between a \emph{new weight query} and a
\emph{new relation}. A reusable unweighted reduction trace avoids a fresh
orbit search, but a new weight query still performs transport through that
trace. On the other hand, adding genuine identifications changes the
relation, so blindly reusing the old final histogram is unsound. We identify
an intermediate interface: new additive observations may be constant on a
fixed source partition, and new relations may collapse whole unions of its
atoms. For this interface, a single compiled table is both sufficient and
closed under all subsequent operations.

Here a \emph{cone} is an equivalence-relation operation, not a statement that
a particular topological cone or manifold attachment is valid. It identifies
every point in a specified union. A gluing that identifies one point in an
interval with one corresponding point in another interval is a different
operation. Keeping these meanings separate is essential to every result
below.

\subsection{Relationship to the inspected repository}

The repository was inspected through its GitHub interface, then pinned at
commit
\nolinkurl{1fab6d945be52857cce48b6459075e6194622001}. The scope is
\path{Topology/UnknotRecognition}; the inspected source paths and blob hashes
are recorded in \path{integration/source_audit.json}. Earlier overview text
lagged behind the newer modules. The present article therefore treats the
newer native sparse and weighted implementations as prior work, rather than
re-presenting their results as new contributions \cite{repo,sparserepo,weightedrepo}.

In particular, the baseline already has a sparse positive component-incidence
interface, a generic weighted orbit histogram, independent weighted replay,
and normal-surface component queries. It also supports reusing a verified
unweighted trace for another weighted query. The contribution here is the
compiled profile interface, its exact quotient action, the sharp height and
compressed-history bounds, and the corresponding additive reference package.
It is not a replacement implementation of the underlying orbit algorithm.

The source was readable, but a runnable native checkout was not available in
the execution environment. Thus the adapter's function signatures were
checked against the pinned files, while its native execution gate remains
\texttt{NOT\_RUN}. The reference implementation and its independent finite
oracles were executed. No maintained production module is modified by this
package, no complete maintained test-suite result is asserted, and no
whole-knot speedup is reported.

\subsection{Classical input and the global target}

Weighted orbit counting is classical. Agol--Hass--Thurston give a polynomial
algorithm in the pairing count, weight-run count, weight dimension, and
binary numerical sizes; their output uses compressed constant-weight blocks,
not a separate expanded record for every represented orbit
\cite[Theorem~16]{aht}. We use this theorem as a compilation primitive.
Union by rank with path compression supplies the standard amortized
set-union bound \cite{tarjan}. Neither ingredient is claimed as a new
algorithmic theorem here.

Lackenby's 2026 hierarchy algorithm is relevant to the larger programme.
Section~9 separates iteration counting from step costs and does not give a
complete runtime estimate. Its Proposition~9.1 bounds iterations using
hierarchy length and pattern complexity; the refinement in Section~10
bounds hierarchy length linearly in initial $q$-complexity
\cite{lackenby}. This does not supply the missing end-to-end
quasi-polynomial implementation bound.

Our narrow question is therefore concrete: after a relation has been encoded
and its future full-port operations described compactly, can the component
observer process the entire phase in time polynomial in that description,
rather than in its expanded number of operations? The answer is yes. The
much stronger question---whether all operations needed by an unknot
recognizer can be organized into a controlled number of such phases---is
not answered here.

\subsection{Main statements at a glance}

The central theorem package has four parts. First, a carry-free scalar weight
recovers the full profile census with binary width determined by the
individual atom capacities. Second, a small graph computes the exact action
of a collection of full-port cones, including multiplicity collapse. Third,
the algebra of these cones has height at most $2m-1$, sharply. Fourth, this
algebra permits direct navigation of compressed operation programs, including
all effective event positions.

The results are presented with proofs, not inferred from experiments. The
experiments instead test the implementation, expose interface mistakes, and
measure the particular kernel improvements. The article makes no assertion
of priority over all literature on finite partition algebras or compressed
program evaluation. Its contribution is the derived theorem package and
implementation for this precisely specified ProveIt observer interface.

\section{Input model, source atoms and bit costs}
\label{sec:model}

Let $X=\{0,\ldots,N-1\}$. An interval pairing has inclusive domain $[a,b]$
and range $[c,d]$, equal lengths, and map either
\[
 x\longmapsto x+c-a
 \qquad\text{or}\qquad
 x\longmapsto a+d-x.
\]
The equivalence relation $E$ is generated by the pairs $(x,f(x))$ for all
supplied pairings $f$. The maps are not expanded point by point. There are
$k$ pairings and every endpoint has at most
$B=\max(1,\lceil\log_2(N+1)\rceil)$ bits.

Fix cuts
\[
 0=c_0<c_1<\cdots<c_m=N,
 \qquad J_j=[c_j,c_{j+1})\cap\Z,
 \qquad \ell_j=c_{j+1}-c_j
 \quad(0\le j<m).
\]
These half-open \emph{source atoms} partition the whole universe. Gaps between
marked ports must be included as atoms; otherwise unmarked components could
have zero profile and be lost. For $N=0$ there are no atoms, no orbits, and a
single cut $0$. This case is handled separately by the software.

Every future port is a union of complete source atoms. If the initial
interface supplies $M$ half-open intervals, their endpoints and the two
universe endpoints yield at most $2M+1$ nonempty atoms. Endpoints of any
additional weight functions must also be inserted before compilation.
Nothing in this construction authorizes a later cut inside an existing atom.

For arithmetic accounting, let $\mathsf M(b)$ bound multiplication of
$b$-bit integers; using schoolbook multiplication gives the safe choice
$\mathsf M(b)=O(b^2)$. Additions and comparisons of $b$-bit counts cost
$O(b)$, not unit time. Set-union operation counts are distinguished from
these integer costs. We use $H$ for the largest bit length of an expanded
program length or a reported event index. A result independent of the
\emph{value} of the expanded length can still depend on $H$.

The implementation is exact Python. Dictionaries give ordinary expected
lookup costs; sorting or balanced dictionaries can replace them with
polynomial comparison overhead. Cooperative cancellation is supported at the
native adapter boundary. The package is not a hostile-input process-memory
sandbox, and none of its count bounds is a timeout guarantee.
